% ------------------------------------------------------------------------
% ------------------------------------------------------------------------
% Modelo UNIAVAN para Artigos Acadêmicos baseado em Modelo SBC
%
% Autor: Luiz Fernando Marquez Arruda
%
% Olá caro autor, segue aqui algumas orientações a respeito de como utilizar este
% template.
%
% Todas as configurações a respeito de nome, curso, orientador e afins devem ser feitas
% no arquivo information.tex que se encontra na pasta setup. Assim, de forma 
% automática, os nomes corretos irão aparecer em sua capa, contra-capa e afins.
%
% Todas as imagens devem submetidas para a pasta images.
%
% A lista de símbolos, siglas e acrônimos também está localizada na pasta setup e
% consta com alguns exemplos para lhe orientar a respeito de como incluir novos
% itens. Obs. os itens já listados podem ser apagados.
%
% De forma a não danificar a estruturação principal, as bibliotecas individuais
% utilizadas podem ser adicionadas no arquivo my_packages.tex localizado na pasta
% setup.
%
% Ao longo do documento você encontrará exemplos de como adicionar uma figura, uma
% tabela, fazer uma citação direta (longa e curta) e indireta.
%
% As listas e sumário são atualizados automaticamente.
% 
% Suas referencias devem ser colocadas no formato bibtex no arquivo references.bib
% na pasta raiz.
%
% Atenciosamente 
%
% Luiz Fernando M. Arruda - 22/04/2022
\documentclass[12pt]{setup/uniavan_article}

% -------------------------------------------------------------------------
% Inclusão de novas bibliotecas
% -------------------------------------------------------------------------
% Inclusão de novas bibliotecas devem ser feitas arquivo setup/my_packages.tex
% -----------------------------

\input{setup/my_packages}

% -------------------------------------------------------------------------
% Dados gerais do documento
% -------------------------------------------------------------------------
% Inclusão do nome do autor, título, orientador, membros da banca, curso etc
% Devem ser editados no arquivo setup/information.tex
% -----------------------------

\input{setup/information}
% ---

% ---
% Configurações de aparência do PDF final
% ---
% alterando o aspecto da cor azul
%\definecolor{blue}{RGB}{41,5,195}
% informações do PDF

% ---

\sloppy

% ----
% Início do documento
% ----
\begin{document}

% Seleciona o idioma do documento (conforme pacotes do babel)
%\selectlanguage{english}
\selectlanguage{brazil}

% Retira espaço extra obsoleto entre as frases.
\frenchspacing 

% Espaçamento 1.5 entre linhas
\onehalfspacing

% -------------------------------------------------------------------------
% ELEMENTOS PRÉ-TEXTUAIS
% -------------------------------------------------------------------------
% Capa, folha de rosto, ficha bibliográfica, errata, folha de aprovação
% Dedicatória, agradecimentos, epígrafe, resumos, listas
% -------------------------------------------------------------------------
% Capa
% -------------------------------------------------------------------------
\imprimircapa
% -------------------------------------------------------------------------


% -------------------------------------------------------------------------
% Folha de rosto
% (o * indica que haverá a ficha bibliográfica)
% -------------------------------------------------------------------------
\imprimirfolhaderosto
% -------------------------------------------------------------------------

\maketitle

% -------------------------------------------------------------------------
% Resumo
% -------------------------------------------------------------------------
\input{beforetext/resumo}
% -------------------------------------------------------------------------


% -------------------------------------------------------------------------
% Abstract
% -------------------------------------------------------------------------
\input{beforetext/abstract}
% -------------------------------------------------------------------------

% ----------------------------------------------------------
% ELEMENTOS TEXTUAIS
% ----------------------------------------------------------

\input{sections/01-section}
\input{sections/02-section}
\input{sections/03-section}
\input{sections/04-section}
\input{sections/05-section}
\input{sections/06-section}
\input{sections/07-section}

% ----------------------------------------------------------
% Referências bibliográficas
% ----------------------------------------------------------
\bibliography{references}

\end{document}
